\hypertarget{appendix-1c7df340573e}{}
\section*{Subtraction through addition}
\section{Subtraction through addition}\label{subtraction-through-addition}

\textbf{Curriculum depth:} Introduced · \textbf{First seen:} C002 · \textbf{Latest development:} C002

\subsection{Definition and mechanism}\label{definition-and-mechanism}

\(A-B=A+\sim B+1\), allowing the same adder to implement both operations; some ISAs fold borrow state into carry-in.

At N-bit residue width, A−B equals A+NOTₙ(B)+1. Binary 6502-style SBC computes A−B−(1−C), equivalently A+NOT₈(B)+C. C=1 means no incoming borrow; output carry=1 means no outgoing borrow.

\subsection{Prerequisites}\label{prerequisites}

\begin{itemize}
\tightlist
\item
  \hyperref[appendix-4fbc03d20829]{Two\textquotesingle s complement} --- For raw unsigned value u in N bits, signed decoding is u when u\textless2ᴺ⁻¹, otherwise u−2ᴺ.
\item
  \hyperref[appendix-d1a5b1e520c2]{Full adder} --- A full adder adds A, B, and carry-in.
\end{itemize}

\subsection{Used by}\label{used-by}

\begin{itemize}
\tightlist
\item
  \hyperref[appendix-89985e6a2685]{Binary long division} --- uses subtraction through addition as part of its representation or reasoning.
\end{itemize}

\subsection{Recall and next depth}\label{recall-and-next-depth}

Explain subtraction through addition without the label, work the example above, then state the assumption whose removal would change the result.

\textbf{Pending:} ISA flag conventions and multi-word subtraction.

\subsection{Source and navigation}\label{source-and-navigation}

\href{Continuous\%20Notes/Computing\%20Foundations\%20-\%20From\%20Bits\%20to\%20C++\%20and\%20Systems\%20APIs\#Arithmetic\%20becomes\%20a\%20policy}{C002 teaching chapter} · \hyperref[appendix-fd090d982a15]{Part definitions and mechanisms} · \href{}{Dependency Map (02-Foundations-II, appendix-4ad47bdd269a)}
